% answers_short.tex —— 2011 年·广东卷·理科 · 简短答案速查
% 由 tools/make_short_answers.py 自 answers.tex 的 \gksolhead/\gkanswer 提取，
% 并与 2011/2011_数学真题_广东理_答案要点.md 逐题比对后生成，勿手改（改源 MD 后重跑）。

\begin{gknotice}
\textbf{使用说明}：本节只给最终答案便于快速核对，完整推导见解析版。答案为编写者独立推导，\textbf{非官方标准答案}。
\end{gknotice}

\gkpart{选择题}
\begin{center}\begin{tabular}{ccccccccc}
\toprule
题号 & 1 & 2 & 3 & 4 & 5 & 6 & 7 & 8 \\
\midrule
答案 & B & C & D & A & C & D & B & A \\
\bottomrule
\end{tabular}\end{center}

\gkpart{填空题}
\begin{itemize}[leftmargin=2.6em]
  \item[\textbf{9.}] $[1,+\infty)$。
  \item[\textbf{10.}] $84$。
  \item[\textbf{11.}] $10$。
  \item[\textbf{12.}] $2$。
  \item[\textbf{13.}] $185$。
  \item[\textbf{14.}] $\left(1,\dfrac{2\sqrt5}{5}\right)$。
\end{itemize}

\gkpart{解答题}
\begin{itemize}[leftmargin=2.6em]
  \item[\textbf{15.}] $\sqrt{35}$。
  \item[\textbf{16.}] (1) $\sqrt2$；(2) $\dfrac{16}{65}$
  \item[\textbf{17.}] (1) $35$ 件；(2) 约 $14$ 件；(3) $P(\xi=0)=\dfrac{3}{10}$，$P(\xi=1)=\dfrac35$，$P(\xi=2)=\dfrac{1}{10}$，$E\xi=\dfrac45$
  \item[\textbf{18.}] (1) 证明见解析；(2) $-\dfrac{\sqrt{21}}{7}$
  \item[\textbf{19.}] (1) $\dfrac{x^2}{4}-y^2=1$；(2) 最大值 $2$，此时 $P\left(\dfrac{6\sqrt5}{5},-\dfrac{2\sqrt5}{5}\right)$
  \item[\textbf{20.}] (1) $b\neq2$ 时 $a_n=\dfrac{nb^n(b-2)}{b^n-2^n}$；$b=2$ 时 $a_n=2$；(2) 证明见解析
  \item[\textbf{21.}] (1) 证明见解析；(2) 证明见解析；(3) $\varphi_{\min}=1$，$\varphi_{\max}=\dfrac54$
\end{itemize}
